\subsection{Ontology Languages: RDF, RDF Schema, OWL}
\label{sec:ontology-languages}
\slidesrc{Ontologies: 24--37}

\subsubsection{Evolution of the languages}
\slidesrc{Ontologies: 25}

Ontology languages for the Web were built on top of each other, starting from the markup
languages of the Web itself.

\begin{center}
\begin{tikzpicture}[
    lay/.style={draw=nbblue!70!black, fill=nbblue!8, minimum height=0.6cm,
                font=\small, inner sep=0pt}]
  \node[lay, minimum width=1.4cm, anchor=south west] at (0,0)    {HTML};
  \node[lay, minimum width=6.6cm, anchor=south west] at (1.4,0)  {XML};
  \node[lay, minimum width=1.6cm, anchor=south west] at (0.7,0.6) {XHTML};
  \node[lay, minimum width=5.4cm, anchor=south west] at (2.6,0.6) {RDF};
  \node[lay, minimum width=4.4cm, anchor=south west] at (3.1,1.2) {RDF Schema};
  \node[lay, minimum width=1.4cm, anchor=south west] at (3.1,1.8) {OWL};
  \node[lay, minimum width=1.4cm, anchor=south west] at (4.6,1.8) {DAML};
  \node[lay, minimum width=1.4cm, anchor=south west] at (6.1,1.8) {OIL};
  \node[lay, minimum width=2.9cm, anchor=south west] at (4.6,2.4) {DAML+OIL};
\end{tikzpicture}
\captionof{figure}{Evolution of the Web and ontology languages: each layer builds on the
ones below it.}
\label{fig:languages-evolution}
\end{center}

\begin{itemize}[nosep]
  \item \textbf{XML} gives a common \emph{syntax} for structured documents, but says
        nothing about what the tags mean.
  \item \textbf{RDF} adds a simple data model (triples, below) on top of XML.
  \item \textbf{RDF Schema} adds a vocabulary for classes, subclasses and properties.
  \item \textbf{DAML} (US) and \textbf{OIL} (Europe) were two richer ontology languages
        built on RDF Schema, merged into \textbf{DAML+OIL}, which in turn was the basis of
        the W3C standard \textbf{OWL}.
\end{itemize}

\subsubsection{RDF: resources, properties, statements}
\slidesrc{Ontologies: 26--29}

\textbf{RDF} (Resource Description Framework) is a W3C Recommendation. It provides a
general, flexible method to decompose any knowledge into small pieces, called
\textbf{triples}, with some rules about the meaning of those pieces.

\begin{nbdef}{RDF data model}
The RDF data model has three object types:
\begin{itemize}[nosep]
  \item a \textbf{resource}: anything that can have a URI (Web identifier). It can be a
        whole Web page, a part of it, a collection of pages, or an object that is not
        directly accessible via the Web (a person, a book);
  \item a \textbf{property}: a specific aspect, characteristic, attribute or relation
        used to describe a resource (\emph{author}, \emph{homepage}). Properties are
        resources too;
  \item a \textbf{statement}: a triple made of a resource, a property and a property
        value (which can be a literal or another resource). The three parts are called
        the \textbf{subject}, \textbf{predicate} and \textbf{object} of the statement.
\end{itemize}
\end{nbdef}

A triple $(O, A, V)$ is more often written as $A(O, V)$, and means: \emph{object $O$ has
attribute $A$ with value $V$}. A set of triples is a directed, labelled \textbf{graph}:
subjects and objects are nodes, each predicate is an edge.

\begin{center}
\begin{tikzpicture}[
    font=\small,
    res/.style={ellipse, draw, align=center, font=\scriptsize\ttfamily, inner sep=1pt},
    lit/.style={rectangle, draw, font=\small, minimum height=0.6cm},
    arr/.style={-{Stealth}, thick},
    lab/.style={font=\scriptsize\ttfamily, fill=white, inner sep=1pt}]
  \node[res] (emp)  at (0,0)    {http://www.w3.org/\\employee/id1321};
  \node[res] (book) at (6.2,0)  {http://www.books.org/\\ISBN0012515866};
  \node[lit] (name) at (2.6,2.0) {Jim Lerners};
  \node[lit] (pr)   at (10.6,0) {\$62};
  \draw[arr] (emp)  -- node[lab, above]{s:authorOf}   (book);
  \draw[arr] (emp)  -- node[lab, left=2pt]{s:hasName} (name);
  \draw[arr] (book) -- node[lab, above]{s:hasPrice}   (pr);
\end{tikzpicture}
\captionof{figure}{An RDF graph with three statements. Ellipses are resources
(identified by URIs), boxes are literal values.}
\label{fig:rdf-graph}
\end{center}

The graph contains three triples:
\begin{align*}
  &\mathit{hasName}(\mathit{id1321}, \text{``Jim Lerners''}),\\
  &\mathit{authorOf}(\mathit{id1321}, \mathit{ISBN0012515866}),\\
  &\mathit{hasPrice}(\mathit{ISBN0012515866}, \text{\$62}).
\end{align*}
Because the book is a
resource, it can be the object of one statement and the subject of another: this is
what makes triples link into a graph.

\begin{nbexample}{RDF in XML syntax}
\begin{minted}{xml}
<?xml version="1.0"?>
<RDF>
  <Description about="http://www.w3schools.com/rdf">
    <author>Jan Egil Refsnes</author>
    <homepage>http://www.w3schools.com</homepage>
  </Description>
</RDF>
\end{minted}
The \texttt{Description} element describes one resource, and each child element is a
property with its value. The document contains two statements:
\begin{center}
\small
\begin{tabular}{@{}lll@{}}
\toprule
\textbf{Subject} & \textbf{Predicate} & \textbf{Object} \\
\midrule
\url{http://www.w3schools.com/rdf} & author   & Jan Egil Refsnes \\
\url{http://www.w3schools.com/rdf} & homepage & \url{http://www.w3schools.com} \\
\bottomrule
\end{tabular}
\end{center}
that is, ``the author of \url{http://www.w3schools.com/rdf} is Jan Egil Refsnes'' and
``the homepage of \url{http://www.w3schools.com/rdf} is
\url{http://www.w3schools.com}''. In the second statement the value is itself a
resource.
\end{nbexample}

\begin{nbwarning}[the example is simplified]
In real RDF/XML the RDF vocabulary lives in its own namespace and the properties in
another one: the root is \texttt{rdf:RDF} with an \texttt{xmlns:rdf} declaration, the
resource is \texttt{rdf:Description rdf:about="\ldots"}, and \texttt{author},
\texttt{homepage} carry a prefix of a vocabulary of our choice. The unprefixed version
above is easier to read but is not valid RDF.
\end{nbwarning}

\subsubsection{RDF Schema}
\slidesrc{Ontologies: 30--31}

RDF alone lets us state facts but not define a vocabulary. \textbf{RDF Schema} (RDFS) is
used to define the structure of the metadata that describe Web resources:
\begin{itemize}[nosep]
  \item it consists of some basic classes and properties, and can be extended to fit
        possibly any domain;
  \item \textbf{classes} are arranged hierarchically (\texttt{rdfs:subClassOf}). The
        root of the hierarchy is \texttt{rdfs:Resource}; \texttt{rdfs:Class} is a
        subclass of \texttt{rdfs:Resource};
  \item \textbf{properties} are instances of the class \texttt{rdf:Property}; they
        describe resources by assigning values to them, and are resources themselves.
        Their use can be constrained to members of certain classes:
        \texttt{rdfs:domain} gives the class of the subjects, \texttt{rdfs:range} the
        class of the values.
\end{itemize}

\begin{nbexample}{teachers and courses in RDFS}
\begin{minted}{xml}
<rdfs:Class rdf:ID="Teacher">
  <rdfs:comment>Teacher Class</rdfs:comment>
  <rdfs:subClassOf rdf:resource="#Person"/>
</rdfs:Class>
<rdfs:Class rdf:ID="Course">
  <rdfs:comment>Course Class</rdfs:comment>
  <rdfs:subClassOf
  rdf:resource="http://www.w3.org/1999/02/22-rdf-syntax-ns#Resource"/>
</rdfs:Class>
<rdf:Property rdf:ID="teacher">
  <rdfs:comment>Teacher of a course</rdfs:comment>
  <rdfs:domain rdf:resource="#Course"/>
  <rdfs:range rdf:resource="#Teacher"/>
</rdf:Property>
<rdf:Property rdf:ID="name">
  <rdfs:comment>Name of a Person or Course</rdfs:comment>
  <rdfs:domain rdf:resource="#Person"/>
  <rdfs:domain rdf:resource="#Course"/>
  <rdfs:range
  rdf:resource="http://www.w3.org/1999/02/22-rdf-syntax-ns#Literal"/>
</rdf:Property>
\end{minted}
\emph{Teacher} is a subclass of \emph{Person}; \emph{teacher} is a property linking a
\emph{Course} (domain) to a \emph{Teacher} (range): ``the teacher of this course is
\ldots''. \emph{name} gives a literal (a string) to persons and courses.
\end{nbexample}

\begin{nbwarning}[two domains mean ``and'', not ``or'']
In RDFS semantics, several \texttt{rdfs:domain} statements on the same property are
\emph{all} applied: whatever has a \emph{name} is inferred to be \emph{both} a
\emph{Person} and a \emph{Course}. To express ``a Person \emph{or} a Course'', the
domain must be a single class that is the union of the two (OWL can define it with
\texttt{owl:unionOf}), or a common superclass.
\end{nbwarning}

\subsubsection{OWL: the Web Ontology Language}
\slidesrc{Ontologies: 32--37}

\textbf{OWL} extends RDF Schema: it uses the RDF meaning of classes and properties
(\texttt{rdfs:Class}, \texttt{rdfs:subClassOf}, \ldots) and adds language primitives for
a richer expressiveness. It comes in three species, which trade expressive power for
ease of reasoning:

\begin{center}
\small
\begin{tabularx}{\linewidth}{@{}l >{\raggedright\arraybackslash}X >{\raggedright\arraybackslash}X >{\raggedright\arraybackslash}X@{}}
\toprule
\textbf{Species} & \textbf{What it is} & \textbf{Advantage} & \textbf{Disadvantage} \\
\midrule
OWL Full & all OWL primitives, combinable in arbitrary ways with RDF and RDF Schema &
  fully upward compatible with RDF & so powerful that it is \textbf{undecidable} \\
OWL DL & a sublanguage of OWL Full that restricts how OWL and RDF constructors can be
  used (DL = description logic) & permits \textbf{efficient reasoning} support & loses
  full compatibility with RDF \\
OWL Lite & a further restriction of OWL DL to a subset of the constructors & easier to
  grasp (users) and to implement (tool builders) & restricted expressivity \\
\bottomrule
\end{tabularx}
\captionof{table}{The three species of OWL, from the most to the least expressive.}
\end{center}

\begin{nbintuition}[why give up expressive power?]
A reasoner answers questions such as ``is class $A$ a subclass of $B$?'' or ``is this
ontology consistent?''. In OWL Full no algorithm can always answer them
(undecidability, as for first-order logic); in OWL DL an algorithm always terminates
with the right answer. Every legal OWL Lite ontology is a legal OWL DL ontology, and
every OWL DL ontology is a legal OWL Full ontology.
\end{nbintuition}

\paragraph{The ontology header.} An OWL document starts by declaring the namespaces it
uses and then describes the ontology itself:
\begin{minted}{xml}
<rdf:RDF
   xmlns:owl ="http://www.w3.org/2002/07/owl#"
   xmlns:rdf ="http://www.w3.org/1999/02/22-rdf-syntax-ns#"
   xmlns:rdfs="http://www.w3.org/2000/01/rdf-schema#"
   xmlns:xsd ="http://www.w3.org/2001/XMLSchema#">
<owl:Ontology rdf:about="">
  <rdfs:comment>An example OWL ontology</rdfs:comment>
  <owl:priorVersion rdf:resource="http://www.mydomain.org/uni-ns-old"/>
  <owl:imports rdf:resource="http://www.mydomain.org/persons"/>
  <rdfs:label>University Ontology</rdfs:label>
</owl:Ontology>
\end{minted}
\texttt{owl:imports} includes other ontologies, whose content is assumed to be part of
the current document. It is \textbf{transitive}: if ontology $A$ imports $B$ and $B$
imports $C$, then $A$ also imports $C$.

\paragraph{Classes.} Classes are defined with \texttt{owl:Class}; further information
can be added to a class later, e.g.\ that it is \textbf{disjoint} from another one (no
individual belongs to both):
\begin{minted}{xml}
<owl:Class rdf:ID="Mammalia">
  <rdfs:subClassOf rdf:resource="#Animal"/>
</owl:Class>
<owl:Class rdf:about="#Mammalia">
  <owl:disjointWith rdf:resource="#Insecta"/>
</owl:Class>
\end{minted}
There are two predefined classes: \texttt{owl:Thing} (the class of everything, the top
of every hierarchy) and \texttt{owl:Nothing} (the empty class).

\paragraph{Properties.} OWL distinguishes two types of property:
\begin{itemize}[nosep]
  \item \textbf{object properties} link individuals to individuals (\emph{married}
        links a person to a person);
  \item \textbf{datatype properties} link individuals to data values (\emph{age} links a
        person to a non-negative integer).
\end{itemize}
\begin{minted}{xml}
<owl:DatatypeProperty rdf:ID="age">
  <rdfs:domain rdf:resource="#Person"/>
  <rdfs:range
  rdf:resource="http://www.w3.org/2001/XMLSchema#nonNegativeInteger"/>
</owl:DatatypeProperty>
<owl:ObjectProperty rdf:ID="married">
  <rdfs:domain rdf:resource="#Person"/>
  <rdfs:range rdf:resource="#Person"/>
</owl:ObjectProperty>
\end{minted}

\begin{nbwarning}[corrections to the slides]
The listings above fix three small slips in the slides: the \texttt{xsd} namespace is
\texttt{XMLSchema} (not \texttt{XLMSchema}); domain and range of OWL properties are
still given with \texttt{rdfs:domain} and \texttt{rdfs:range} (OWL has no
\texttt{owl:domain}/\texttt{owl:range}); and \texttt{rdf:about} on a class defined
elsewhere refers to it with a \texttt{\#}.
\end{nbwarning}

\paragraph{References.} The specifications are on the W3C site: RDF
(\url{http://www.w3.org/RDF/}), RDF Schema (\url{http://www.w3.org/TR/rdf-schema/}) and
OWL (\url{http://www.w3.org/2004/OWL/}).
